\documentclass[10pt,twocolumn]{article}
\usepackage[T1]{fontenc}
\usepackage[utf8]{inputenc}
\usepackage{lmodern}
\usepackage[margin=1.8cm,top=2.4cm,bottom=2cm,columnsep=0.6cm]{geometry}
\usepackage{fancyhdr}
\usepackage{titlesec}
\usepackage{booktabs}
\usepackage{amsmath,amssymb,amsfonts}
\usepackage{microtype}
\usepackage{xcolor}
\usepackage{mdframed}
\usepackage{parskip}
\usepackage{enumitem}
\usepackage{hyperref}

\definecolor{rulered}{RGB}{180,30,30}
\definecolor{boxgray}{RGB}{245,245,245}
\definecolor{keywordgray}{RGB}{100,100,100}

\pagestyle{fancy}
\fancyhf{}
\fancyhead[L]{\textsc{\small Claude Mag}}
\fancyhead[C]{\small\textit{Machine Learning Research Digest}}
\fancyhead[R]{\small 2026-10-05}
\fancyfoot[C]{\small\thepage}
\renewcommand{\headrulewidth}{0.8pt}

\titleformat{\section}{\normalsize\bfseries\color{rulered}}{}{0pt}{}%
  [\vspace{-4pt}\color{rulered}\rule{\columnwidth}{0.4pt}]
\titlespacing{\section}{0pt}{8pt}{4pt}

\hypersetup{colorlinks=true,urlcolor=rulered,linkcolor=black}

\begin{document}
\setlength{\parindent}{0pt}
\setlength{\parskip}{4pt}

{\small\color{keywordgray}\textbf{Keywords:}
  memory-efficient fine-tuning \textbullet{}
  operator-norm optimization \textbullet{}
  ternary update \textbullet{}
  LLM training}\par\vspace{4pt}

{\LARGE\bfseries\raggedright
  TACO: Ternary Absolute-max Column-wise One-sparse Optimizer for LLM Fine-Tuning}\par\vspace{4pt}

{\large\itshape\raggedright
  174$\times$ Less Optimizer State Than AdamW8bit, Full-Parameter Fine-Tuning of 32B Models
  on One 80\,GB H100, Comparable Accuracy}\par\vspace{6pt}

{\small\textbf{By} Jichao Jiang, Cristian McGee, El Houcine Bergou,
  Hanqin Cai \& Aritra Dutta (University of Central Florida;
  Mohammed VI Polytechnic University)}\par\vspace{2pt}

{\small\color{keywordgray}arXiv:2610.02199 \textbullet{} October 2026}
\par\vspace{2pt}\noindent{\color{rulered}\rule{\columnwidth}{0.6pt}}\vspace{6pt}

TACO takes the operator-norm steepest-descent view to its logical extreme:
it selects the sign of the largest-magnitude entry in each column of every
two-dimensional weight matrix---a ternary, one-sparse-per-column update that
requires essentially zero persistent optimizer state---and reduces AdamW8bit's
27.7\,GB optimizer footprint on OPT-13B to 0.16\,GB while matching downstream
accuracy across multiple model families and enabling 32B-scale full-parameter
fine-tuning on a single H100.

\begin{mdframed}[backgroundcolor=boxgray,linecolor=rulered,linewidth=1pt,
    innerleftmargin=6pt,innerrightmargin=6pt,
    innertopmargin=6pt,innerbottommargin=6pt]
\textbf{\small AT A GLANCE}\par\vspace{4pt}
\begin{description}[leftmargin=0pt,labelsep=4pt,itemsep=2pt,parsep=0pt,
    font=\small\bfseries]
  \item[Problem:] AdamW and quantized variants require 2--3$\times$ the model size
    in optimizer state, limiting full-parameter fine-tuning to small models on
    a single GPU.
  \item[Why it matters:] Full-parameter fine-tuning outperforms LoRA on complex
    domain adaptation; enabling it for 30B+ models democratizes access.
  \item[Method:] Exact steepest descent under dimension-normalized 1$\to$1 operator
    norm selects sign of the maximum-magnitude column entry; no persistent momentum
    buffers.
  \item[Result:] 174$\times$ optimizer state reduction; 2.9$\times$ peak memory
    reduction; accuracy within 0.5--1.0\% of AdamW on SuperGLUE, SQuAD v2.
  \item[Evidence:] OPT, LLaMA, Mistral, Phi families; fine-tuning tasks including
    SuperGLUE, SQuAD v2, and instruction following.
\end{description}
\end{mdframed}

\section{The Big Picture}

Full-parameter fine-tuning of large language models is the gold standard for
domain adaptation: it updates every weight, avoids the rank constraints of
parameter-efficient methods (LoRA, prefix tuning), and captures distribution
shifts that low-rank approximations cannot represent. Yet it has been practically
out of reach for models above $\sim$7B parameters on a single 80\,GB GPU because
AdamW---the dominant optimizer---maintains first and second moment buffers that
add 2$\times$ the model parameter count in optimizer state.

TACO addresses this by rethinking what state an optimizer actually needs.
Drawing on the operator-norm steepest-descent perspective introduced by Muon,
TACO shows that the exact optimal update under a dimension-normalized 1$\to$1
operator norm requires no momentum buffers whatsoever---only the current-step
gradient, which is discarded immediately after use.

\section{Why Existing Approaches Fall Short}

\textbf{AdamW:} Two full-precision moment buffers (first and second order) add
$2\times$ the model's FP32 parameter count to GPU memory. On OPT-13B this is
27.7\,GB of persistent state alone.

\textbf{AdamW8bit:} Quantizes moment buffers to 8 bits, cutting state to
$\approx 27.7$\,GB/$2$\,=\,13.9\,GB---a useful reduction but still far larger
than TACO's 0.16\,GB.

\textbf{SignSGD:} Takes the sign of the full gradient column-by-column but does
not arise from a principled operator-norm steepest-descent argument, and can
exhibit slow convergence on poorly scaled gradients.

\textbf{Muon:} Derives updates via operator-norm steepest descent and applies
Nesterov momentum, but still maintains a first-moment buffer.

\textbf{LoRA:} Reduces state by fine-tuning only low-rank adaptors, but cannot
update all parameters and underperforms full-rank methods on complex adaptation.

\section{Core Mechanism}

\textbf{Operator-norm steepest descent.} Given a weight matrix $W \in \mathbb{R}^{m\times n}$
with gradient $G$, the steepest-descent direction under the dimension-normalized
$1 \to 1$ operator norm constraint selects, for each column $j$, the row index
achieving the maximum gradient magnitude, and assigns a ternary entry equal to
the sign of that element.

Formally, the update matrix $U \in \{-1, 0, +1\}^{m \times n}$ has exactly
one nonzero entry per column---hence ``one-sparse per column'' and ``ternary''.
No previous-step information is accumulated: the only state at any step is the
current gradient, discarded immediately after $U$ is computed.

\section{Technical Formulation}

For each column $j \in [n]$ of gradient $G^{(t)}$:
\[
i^*(j) = \arg\max_{i \in [m]} |G_{ij}^{(t)}|
\]
\[
U_{ij}^{(t)} = \begin{cases}
  \mathrm{sign}\!\left(G_{i^*(j),j}^{(t)}\right) & i = i^*(j) \\
  0 & \text{otherwise}
\end{cases}
\]
Weight update:
\[
W^{(t+1)} = W^{(t)} - \eta\, U^{(t)}
\]
Persistent optimizer state between steps: $\varnothing$ (the BF16 gradient
buffer $G^{(t)}$ is freed immediately after the update).

Under a convex loss with Lipschitz gradient, TACO achieves
$O(1/\sqrt{T})$ convergence matching signSGD, while the update $U^{(t)}$
is the exact minimizer of the linearized objective subject to the unit
operator-norm ball.

\section{Learning or Inference Procedure}

\textbf{Forward:} Standard.

\textbf{Backward:} Compute gradient $G^{(t)}$ for each 2D weight matrix.

\textbf{Update (per-matrix):}
\begin{enumerate}[itemsep=2pt,parsep=0pt]
  \item Compute argmax over each column of $|G^{(t)}|$.
  \item Construct ternary one-sparse matrix $U^{(t)}$.
  \item Apply $W \leftarrow W - \eta U^{(t)}$ in-place.
  \item Free $G^{(t)}$.
\end{enumerate}

Non-2D tensors (embeddings, 1D biases) fall back to signSGD.
Learning rate: cosine or linear warmup-decay schedule; no additional
hyperparameters beyond $\eta$.

\section{What the Guarantee Says}

For convex losses, TACO converges at rate $O(1/\sqrt{T})$, matching signSGD.
The $U^{(t)}$ is the exact steepest-descent direction under the 1$\to$1
operator norm: no other ternary update achieves a greater linearized objective
decrease under the operator-norm geometry. The persistent state between steps
is strictly zero beyond the model weights, providing a hard lower bound on
optimizer memory that no further quantization can improve.

\section{Experimental Findings}

\begin{itemize}[itemsep=2pt,parsep=0pt]
  \item \textbf{Optimizer state:} 27.7\,GB (AdamW8bit) $\to$ 0.16\,GB (TACO)
    on OPT-13B: a 174$\times$ reduction.
  \item \textbf{Peak training memory:} 80.6\,GB $\to$ 27.5\,GB on OPT-13B:
    a 2.9$\times$ reduction.
  \item \textbf{Large-model access:} OPT-30B and LLaMA-32B can be fully
    fine-tuned on a single 80\,GB H100---previously impossible with AdamW.
  \item \textbf{Accuracy:} Within 0.5--1.0\% of AdamW on SuperGLUE, SQuAD\,v2,
    and instruction following, across OPT, LLaMA, Mistral, Phi.
  \item \textbf{Throughput:} Training step runtime within 5\% of AdamW8bit
    owing to the simplicity of the argmax-and-sign update.
\end{itemize}

\section{Ablations and Interpretation}

\textbf{Row-wise vs.\ column-wise argmax:} Row-wise argmax (violating the
operator-norm derivation) slightly degrades accuracy, confirming the
theoretical motivation matters in practice.

\textbf{Top-$k$ per column (TACO-$k$):} Allowing $k>1$ nonzero entries per
column interpolates between TACO and a dense gradient, recovering AdamW accuracy
at the cost of $k$-fold increased state; $k=1$ is Pareto-optimal for
memory-constrained settings.

\textbf{Momentum ablation:} Adding a first-moment moving average of $U^{(t)}$
improves convergence speed on some tasks but doubles persistent state; omitting
momentum is preferred when memory is the bottleneck.

\textbf{TACO vs.\ LoRA:} At equal GPU memory budget, full-parameter TACO
outperforms LoRA on domain-specific tasks, confirming the value of full-rank updates
when memory constraints are relaxed by the optimizer.

\section*{Reference}

Jichao Jiang, Cristian McGee, El Houcine Bergou, Hanqin Cai, Aritra Dutta.
\textbf{TACO: Ternary Absolute-max Column-wise One-sparse Optimizer for LLM Fine-Tuning}.
arXiv:2610.02199, October 2026.\\
\url{https://arxiv.org/abs/2610.02199}

\end{document}
